\documentclass[12pt,a4paper]{article}
\usepackage[utf8]{inputenc}
\usepackage[T1]{fontenc}
\usepackage[spanish]{babel}
\usepackage{amsmath,amssymb}
\usepackage{geometry}
\usepackage{booktabs}
\usepackage{array}
\usepackage{longtable}
\usepackage{hyperref}
\geometry{margin=2.5cm}

\title{Guia Completa de Defensa Matematica\\
Actividad 1 -- Simulacion por Transformada Inversa}
\author{Basada en Teorema.tex e Informe.tex}
\date{\today}

\begin{document}
\maketitle
\tableofcontents
\newpage

\section{Objetivo de esta guia}
Esta guia resume \textbf{todo el sustento matematico y de simulacion} de la actividad, con enfoque de defensa oral en clase. Incluye:
\begin{itemize}
  \item Teorema utilizado y su interpretacion.
  \item Desarrollo paso a paso de la Parte 1.
  \item Desarrollo matematico y de costos de Parte 2, Problema 1.
  \item Desarrollo matematico y de costos de Parte 2, Problema 2.
  \item Significado de cada variable y respuesta corta para preguntas del docente.
\end{itemize}

\section{Base teorica: metodo de la transformada inversa}

\subsection{Idea central del teorema}
El metodo de la transformada inversa (capitulo de Coss Bu) establece que, si una variable aleatoria continua $X$ tiene distribucion acumulada $F(x)$ y si $R\sim U(0,1)$, entonces:
\[
X = F^{-1}(R)
\]
genera valores con la distribucion de $X$.

\subsection{Pasos generales del metodo}
\begin{enumerate}
  \item Definir la densidad $f(x)$ y su dominio.
  \item Calcular la acumulada: $F(x)=\int f(x)\,dx$ en el dominio correcto.
  \item Igualar $F(x)=R$, con $R\in[0,1]$ uniforme.
  \item Despejar $x$ para obtener $F^{-1}(R)$.
  \item Verificar extremos ($R=0$ y $R=1$) para validar el intervalo de salida.
\end{enumerate}

\subsection{Ejemplos de Teorema.tex que sustentan la actividad}
\begin{itemize}
  \item \textbf{Exponencial:} 
  \[
  f(x)=\lambda e^{-\lambda x},\ x\ge 0,\quad x=-\frac{1}{\lambda}\ln(R)
  \]
  (Usado directamente en Problema 2 para tiempos entre llegadas).

  \item \textbf{Uniforme en intervalo $[a,b]$:}
  \[
  x=a+(b-a)R
  \]
  (Usado en Problema 1 y en tiempos de descarga del Problema 2).

  \item \textbf{Distribucion por tramos/discreta:}
  se construye acumulada y se asigna el valor segun el rango de $R$.
\end{itemize}

\section{Parte 1: transformada inversa para la densidad dada}

\subsection{Definicion de la densidad}
La funcion objetivo es:
\[
f(x)=
\begin{cases}
\dfrac{(x-3)^2}{18}, & 0\le x\le 6,\\
0, & \text{en otro caso.}
\end{cases}
\]

\subsection{Validacion de densidad}
Para que sea una densidad valida debe cumplirse:
\[
\int_0^6 f(x)\,dx = 1
\]
Desarrollo:
\[
\int_0^6 \frac{(x-3)^2}{18}\,dx
=\frac{1}{18}\left[\frac{(x-3)^3}{3}\right]_0^6
=\frac{1}{18}\left(\frac{27}{3}-\frac{-27}{3}\right)
=\frac{18}{18}=1.
\]
Conclusiones:
\begin{itemize}
  \item $f(x)\ge 0$ en $[0,6]$.
  \item El area total es 1.
  \item Entonces $f(x)$ es una funcion de densidad valida.
\end{itemize}

\subsection{Acumulada $F(x)$ paso a paso}
\[
F(x)=\int_0^x \frac{(t-3)^2}{18}\,dt
\]
\begin{align*}
F(x)
&=\frac{1}{18}\int_0^x (t-3)^2\,dt \\
&=\frac{1}{18}\left[\frac{(t-3)^3}{3}\right]_0^x \\
&=\frac{1}{18}\left(\frac{(x-3)^3}{3}-\frac{(0-3)^3}{3}\right) \\
&=\frac{1}{18}\left(\frac{(x-3)^3}{3}+\frac{27}{3}\right) \\
&=\frac{(x-3)^3+27}{54}.
\end{align*}
Por tanto:
\[
\boxed{F(x)=\frac{(x-3)^3+27}{54}},\quad 0\le x\le 6.
\]

\subsection{Inversa $F^{-1}(R)$ paso a paso}
\[
R=\frac{(x-3)^3+27}{54}
\]
\begin{align*}
54R&=(x-3)^3+27 \\
(x-3)^3&=54R-27 \\
x-3&=\sqrt[3]{54R-27} \\
\boxed{x=3+\sqrt[3]{54R-27}}
\end{align*}
Esta formula es la que se programa para generar los valores de la Parte 1.

\subsection{Verificacion del intervalo de salida}
\begin{itemize}
  \item Si $R=0$: $x=3+\sqrt[3]{-27}=0$.
  \item Si $R=1$: $x=3+\sqrt[3]{27}=6$.
\end{itemize}
Como $R\in[0,1]$, entonces $x\in[0,6]$.

\subsection{Valor esperado teorico}
\[
E[X]=\int_0^6 x\cdot\frac{(x-3)^2}{18}\,dx
\]
Expandiendo $(x-3)^2=x^2-6x+9$:
\[
E[X]=\frac{1}{18}\int_0^6 (x^3-6x^2+9x)\,dx
\]
\[
=\frac{1}{18}\left[\frac{x^4}{4}-2x^3+\frac{9x^2}{2}\right]_0^6
=\frac{1}{18}(324-432+162)=\frac{54}{18}=3.
\]
\[
\boxed{\mu=E[X]=3}
\]

\subsection{Interpretacion de la validacion por simulacion}
La media muestral $\bar x$ obtenida para $n=10,25,50,100,250,500,1000$ debe acercarse a 3 al crecer $n$ (Ley de los Grandes Numeros). La prueba de escritorio documenta esa convergencia.

\section{Parte 2, Problema 1: maquinas y mecanicos}

\subsection{Enunciado matematico del sistema}
Se evalua el costo total por turno de 8 horas segun cuantas maquinas ($m$) atiende un mecanico.

Datos de costo:
\begin{itemize}
  \item Costo de maquina ociosa: \$500 por hora.
  \item Salario del mecanico: \$50 por hora.
  \item Turno: 8 horas $\Rightarrow$ salario fijo por turno = $8\times 50=\$400$.
\end{itemize}

\subsection{Variables que debes saber explicar}
\begin{itemize}
  \item $R_1$: uniforme para seleccionar tiempo entre descomposturas.
  \item $t_d$: tiempo entre descomposturas.
  \item $R_2$: uniforme para seleccionar tiempo de reparacion.
  \item $t_r$: tiempo de reparacion.
  \item $H_{\text{inact}}$: horas ociosas acumuladas.
  \item $m$: numero de maquinas por mecanico.
\end{itemize}

\subsection{Modelo de generacion de tiempos}
Para cada intervalo $[a,b]$ asignado por la acumulada, se usa uniforme:
\[
t=a+(b-a)R
\]
Esto aplica tanto para $t_d$ como para $t_r$, cambiando el intervalo segun la tabla de probabilidades.

\subsection{Funcion de costo del problema 1}
El costo total por turno para un escenario $m$ es:
\[
\text{Costo total}(m)=\underbrace{500\cdot H_{\text{inact}}(m)}_{\text{costo por inactividad}}+\underbrace{400}_{\text{salario fijo del mecanico}}.
\]

\subsection{Interpretacion economica}
\begin{itemize}
  \item Si aumenta $m$, el mecanico atiende mas carga y puede crecer inactividad.
  \item El salario fijo no cambia (siempre un mecanico).
  \item El escenario optimo es el de menor costo total promedio simulado.
\end{itemize}

\subsection{Como defender la prueba de escritorio}
En tus tablas:
\begin{itemize}
  \item Para $n=10$ y $n=25$: se muestran todas las replicas.
  \item Para $n\ge 50$: se resume (corridas iniciales + ultima + promedio).
\end{itemize}
Argumento clave: con pocas replicas hay alta variabilidad por eventos extremos; con $n$ grande se estabiliza el promedio, lo que permite tomar decision robusta.

\section{Parte 2, Problema 2: equipo de trabajadores en almacen}

\subsection{Enunciado matematico}
Llegan camiones con proceso Poisson de tasa:
\[
\lambda=2\ \text{camiones/hora}
\]
Se comparan escenarios de $w\in\{3,4,5,6\}$ trabajadores para minimizar costo total por turno de 8 horas.

\subsection{Variables que debes saber explicar}
\begin{itemize}
  \item $\lambda$: tasa de llegadas del proceso Poisson.
  \item $R$: uniforme para transformar a tiempo entre llegadas.
  \item $t_{\text{llegada}}$: tiempo entre llegadas de camiones.
  \item $w$: numero de trabajadores.
  \item $T_{\text{espera}}$: tiempo total de espera acumulado de camiones.
\end{itemize}

\subsection{Transformada usada (Teorema.tex, ejemplo exponencial)}
Si llegadas son Poisson, los tiempos entre llegadas son exponenciales:
\[
t_{\text{llegada}}=-\frac{1}{\lambda}\ln(R)=-\frac{1}{2}\ln(R).
\]
Esto se obtiene por transformada inversa de la exponencial.

\subsection{Tiempos de descarga por escenario}
Se modelan con uniformes:
\begin{itemize}
  \item $w=3$: $t=20+10R$ (min).
  \item $w=4$: $t=15+10R$ (min).
  \item $w=5$: $t=10+10R$ (min).
  \item $w=6$: $t=5+10R$ (min).
\end{itemize}

\subsection{Funcion de costo del problema 2}
Con salario de \$25/h por trabajador y costo de espera de \$50/h:
\[
\text{Costo total}(w)=\underbrace{200w}_{\text{salarios en 8h}}+\underbrace{50\,T_{\text{espera}}(w)}_{\text{espera de camiones}}.
\]

\subsection{Interpretacion economica}
\begin{itemize}
  \item Si $w$ es bajo, sube espera (colas), baja salario.
  \item Si $w$ es alto, baja espera, pero sube salario.
  \item El optimo aparece en el equilibrio entre ambos componentes.
\end{itemize}

\section{Relacion exacta entre teoria y codigo}

\subsection{Parte 1}
Codigo implementa:
\[
x=3+\sqrt[3]{54R-27}
\]
Cada corrida toma un $R$ uniforme y transforma a $x$.

\subsection{Problema 1}
Codigo implementa:
\begin{enumerate}
  \item Seleccion de intervalos via acumulada para descompostura y reparacion.
  \item Generacion de tiempos via $t=a+(b-a)R$.
  \item Acumulacion de inactividad y costo.
  \item Comparacion de escenarios por promedio.
\end{enumerate}

\subsection{Problema 2}
Codigo implementa:
\begin{enumerate}
  \item Llegadas via $t=-\ln(R)/\lambda$.
  \item Descargas via uniformes por escenario de $w$.
  \item Calculo de espera por cola (FIFO) y costo total.
  \item Comparacion de escenarios por promedio.
\end{enumerate}

\section{Preguntas de defensa (respuestas cortas)}

\subsection*{Pregunta 1: por que funciona $F^{-1}(R)$?}
Porque transforma una uniforme en la distribucion objetivo usando la acumulada de esa variable.

\subsection*{Pregunta 2: por que verificas integral igual a 1?}
Porque es condicion necesaria para que $f(x)$ sea una densidad valida.

\subsection*{Pregunta 3: por que la media teorica en Parte 1 es 3?}
Porque al integrar $x f(x)$ en $[0,6]$ el resultado exacto es $E[X]=3$.

\subsection*{Pregunta 4: por que no decidir con una sola corrida?}
Porque una corrida aislada puede tener alta variabilidad; el promedio de muchas replicas reduce sesgo por azar.

\subsection*{Pregunta 5: que representa $\lambda=2$?}
Tasa media de 2 camiones por hora; su inversa es 0.5 horas entre llegadas en promedio.

\subsection*{Pregunta 6: por que usar uniforme por intervalos en Problema 1?}
Porque el enunciado da rangos con probabilidades; primero se escoge intervalo con acumulada y luego valor interno con uniforme.

\section{Checklist de exposicion oral}
\begin{enumerate}
  \item Explicar el teorema en 20 segundos.
  \item Derivar en pizarra $F(x)$ y $F^{-1}(R)$ de Parte 1.
  \item Verificar extremos $R=0$ y $R=1$.
  \item Mostrar que $E[X]=3$ teorico.
  \item Explicar funciones de costo de Problema 1 y Problema 2.
  \item Explicar por que se comparan promedios con multiples replicas.
  \item Cerrar con criterio de decision: minimizar costo total esperado.
\end{enumerate}

\section{Conclusion tecnica}
La actividad aplica correctamente el marco de Coss Bu:
\begin{itemize}
  \item Transformada inversa para distribuciones continuas.
  \item Acumuladas para seleccion por intervalos.
  \item Simulacion Monte Carlo para estimar costos esperados.
\end{itemize}
En consecuencia, cada decision (escenario optimo) se fundamenta en un modelo probabilistico, una funcion de costo y evidencia por replicas suficientes.

\end{document}
